\documentclass[12pt,a4paper]{article}

\usepackage[margin=1in]{geometry}
\usepackage{amsmath,amssymb,amsthm}
\usepackage{enumitem}
\usepackage{hyperref}

\hypersetup{
  colorlinks=true,
  linkcolor=blue,
  urlcolor=blue,
  citecolor=blue
}

\newtheorem{theorem}{Theorem}

\theoremstyle{remark}
\newtheorem*{idea}{Key Idea}

\title{Maschke's Theorem and Its Proof}
\author{}
\date{}

\begin{document}

\maketitle

\begin{theorem}[Maschke]
Let \(G\) be a finite group, and let
\[
  \rho: G \longrightarrow \operatorname{GL}(V)
\]
be a finite-dimensional representation of \(G\) over a field \(\mathbb{F}\).
If
\[
  \operatorname{char}(\mathbb{F}) \nmid |G|,
\]
then \(\rho\) is completely reducible. More precisely, for every
\(G\)-invariant subspace \(W \subseteq V\), there exists a \(G\)-invariant
subspace \(U \subseteq V\) such that
\[
  V = W \oplus U.
\]
Equivalently, every \(G\)-invariant subspace has a \(G\)-invariant
complement.
\end{theorem}

\begin{proof}
The central idea of the proof is \emph{averaging}. Since the group is finite,
we can average an arbitrary linear projection over the group action and thereby
turn it into a \(G\)-equivariant projection.

\subsection*{Step 1: Construct an arbitrary linear projection}

Let \(W \subseteq V\) be a \(G\)-invariant subspace. By linear algebra, we may
choose a basis
\[
  \{w_1,\dots,w_m\}
\]
of \(W\), and extend it to a basis
\[
  \{w_1,\dots,w_m,u_1,\dots,u_k\}
\]
of \(V\). Define a linear map
\[
  \pi: V \longrightarrow W
\]
by
\[
  \pi(w_i)=w_i,\qquad \pi(u_j)=0,
  \qquad i=1,\dots,m,\ j=1,\dots,k.
\]
Then
\[
  \pi|_W=\operatorname{id}_W,\qquad
  \operatorname{im}(\pi)=W,\qquad
  \pi^2=\pi.
\]
Thus \(\pi\) is a projection onto \(W\). However, \(\pi\) is not necessarily
\(G\)-equivariant; it is merely a linear projection.

\subsection*{Step 2: Average the projection}

Define a new linear map \(\widetilde{\pi}:V\to V\) by
\[
  \widetilde{\pi}(v)
  =
  \frac{1}{|G|}
  \sum_{g\in G}
  \rho(g)\,\pi\bigl(\rho(g)^{-1}v\bigr),
  \qquad v\in V.
\]
The scalar \(|G|^{-1}\) exists in \(\mathbb{F}\), because
\(\operatorname{char}(\mathbb{F})\nmid |G|\).

Intuitively, for each \(g\in G\), we first pull \(v\) back by
\(\rho(g)^{-1}\), project it onto \(W\) using \(\pi\), push the result forward
by \(\rho(g)\), and then average over all group elements.

\subsection*{Step 3: Verify the key properties of \(\widetilde{\pi}\)}

\paragraph{(a) The image of \(\widetilde{\pi}\) is contained in \(W\).}
For any \(v\in V\) and \(g\in G\), we have
\[
  \pi\bigl(\rho(g)^{-1}v\bigr)\in W.
\]
Since \(W\) is \(G\)-invariant, applying \(\rho(g)\) keeps the vector inside
\(W\). Hence
\[
  \rho(g)\,\pi\bigl(\rho(g)^{-1}v\bigr)\in W.
\]
Therefore the finite sum also lies in \(W\), and so
\[
  \operatorname{im}(\widetilde{\pi})\subseteq W.
\]

\paragraph{(b) The map \(\widetilde{\pi}\) restricts to the identity on \(W\).}
Let \(w\in W\). Since \(W\) is \(G\)-invariant, we have
\[
  \rho(g)^{-1}w\in W.
\]
Because \(\pi|_W=\operatorname{id}_W\),
\[
  \pi\bigl(\rho(g)^{-1}w\bigr)=\rho(g)^{-1}w.
\]
Therefore
\[
  \rho(g)\,\pi\bigl(\rho(g)^{-1}w\bigr)
  =
  \rho(g)\rho(g)^{-1}w
  =
  w.
\]
Thus every term in the average equals \(w\), so
\[
  \widetilde{\pi}(w)
  =
  \frac{1}{|G|}
  \sum_{g\in G} w
  =
  w.
\]
Consequently,
\[
  \widetilde{\pi}|_W=\operatorname{id}_W.
\]
Together with \(\operatorname{im}(\widetilde{\pi})\subseteq W\), this implies
\[
  \operatorname{im}(\widetilde{\pi})=W
  \qquad\text{and}\qquad
  \widetilde{\pi}^2=\widetilde{\pi}.
\]

\paragraph{(c) The map \(\widetilde{\pi}\) is \(G\)-equivariant.}
Let \(h\in G\) and \(v\in V\). Then
\[
\begin{aligned}
  \widetilde{\pi}(\rho(h)v)
  &=
  \frac{1}{|G|}
  \sum_{g\in G}
  \rho(g)\,\pi\bigl(\rho(g)^{-1}\rho(h)v\bigr) \\
  &=
  \frac{1}{|G|}
  \sum_{g\in G}
  \rho(g)\,\pi\bigl(\rho(g^{-1}h)v\bigr).
\end{aligned}
\]
Now set \(t=h^{-1}g\), so \(g=ht\). As \(g\) ranges over \(G\), so does
\(t\). Hence
\[
\begin{aligned}
  \widetilde{\pi}(\rho(h)v)
  &=
  \frac{1}{|G|}
  \sum_{t\in G}
  \rho(ht)\,\pi\bigl(\rho(t)^{-1}v\bigr) \\
  &=
  \rho(h)
  \left(
  \frac{1}{|G|}
  \sum_{t\in G}
  \rho(t)\,\pi\bigl(\rho(t)^{-1}v\bigr)
  \right) \\
  &=
  \rho(h)\,\widetilde{\pi}(v).
\end{aligned}
\]
Therefore \(\widetilde{\pi}\) is a \(G\)-equivariant projection onto \(W\).

\subsection*{Step 4: Take the kernel to obtain a \(G\)-invariant complement}

Define
\[
  U=\ker(\widetilde{\pi})
  =
  \{u\in V\mid \widetilde{\pi}(u)=0\}.
\]
Since \(\widetilde{\pi}\) is a projection with image \(W\), linear algebra gives
\[
  V=\operatorname{im}(\widetilde{\pi})\oplus\ker(\widetilde{\pi})
  =
  W\oplus U.
\]
It remains only to check that \(U\) is \(G\)-invariant. If \(u\in U\), then
\(\widetilde{\pi}(u)=0\). For any \(g\in G\), using \(G\)-equivariance gives
\[
  \widetilde{\pi}(\rho(g)u)
  =
  \rho(g)\widetilde{\pi}(u)
  =
  \rho(g)0
  =
  0.
\]
Thus \(\rho(g)u\in U\), so \(U\) is \(G\)-invariant.

We have therefore found a \(G\)-invariant complement \(U\) of \(W\), namely
\[
  V=W\oplus U.
\]
\end{proof}

\begin{idea}
The proof rests on four points:
\begin{enumerate}[label=\arabic*.]
  \item Linear algebra provides an arbitrary projection \(\pi:V\to W\).
  \item Averaging over the finite group turns \(\pi\) into a \(G\)-equivariant
  projection \(\widetilde{\pi}\).
  \item The change of variables in the group sum proves equivariance.
  \item The kernel of \(\widetilde{\pi}\) is automatically a \(G\)-invariant
  complement of \(W\).
\end{enumerate}
\end{idea}

\end{document}
